\chapter*{原书序}
\addcontentsline{toc}{chapter}{原书序}
\markboth{原书序}{原书序}

过去半个世纪里，自守形式理论已成为数论与代数几何发展的一个主要焦点，其应用遍及许多不同领域，包括组合学与数学物理。

本书所收 12 章，基于 2001 年 3 月 12--16 日在耶路撒冷希伯来大学高等研究院（Institute for Advanced Studies）数学学院所作的系列讲座。这些讲座的目标，是向青年研究者介绍自守形式理论，说明它与 L 函数的联系，并指出它与其他数学领域的关联。

本次讲习班的核心主旨是阐述朗兰兹纲领——它给出了一幅极为宏阔的图景，把自守形式与 L 函数联系起来。为描述这幅图景，各位讲者动用了不同的技术手段，来建立该纲领的一些特殊情形。

朗兰兹纲领粗略地说：在诸多内容之中，任何用数论方式定义的 L 函数，都等同于某个 $\GL(n)$ 的自守 L 函数。以这种宽松的方式看，每个 L 函数（在猜想的意义上）都可视为同一个对象。

为介绍自守形式理论，我们并未试图给出该主题完整而正式的导引；那将需要多得多的篇幅。讲者们转而聚焦于一系列专题，以「非正式」的个人视角，从经典的 zeta 函数一路讲到朗兰兹纲领。我们希望本书也能起到同样非正式的作用。

各位讲者的贡献与其文章，可简述如下。

前八章处理 $\GL(1)$ 与 $\GL(2)$ 的情形：E.~Kowalski 沿经典路线，聚焦 Riemann 的基本 zeta 函数及其向 Dirichlet L 函数与 Hecke L 函数的推广、类域论，以及若干关于经典自守函数的专题；E.~de Shalit 细致综述了 Artin 猜想与 Shimura--Taniyama--Weil 猜想。

在讨论 Hecke L 函数之后，S.~Kudla 把经典模（自守）L 函数作为 $\GL(2)$ 的 L 函数来考察，由此引入表示论。研究那些「自守」表示的一条途径，是 Selberg 的「迹公式」理论；这由 D.~Bump 的一章引入。

最后四章由 J.~Cogdell 与 D.~Gaitsgory 撰写，更为抽象。Cogdell 从讨论 $\GL(2,\A)$ 的尖点自守表示入手，很快推进到 $\GL(n,\A)$ 的朗兰兹理论；随后他解释了为什么需要朗兰兹对偶群，才能对不同于 $\GL(n)$ 的约化群 $G$ 陈述一般猜想。

Gaitsgory 对几何朗兰兹纲领作了非正式的介绍。这是一个新兴而活跃的研究领域，它由自守形式理论生长出来，并与后者密切相关。粗略地说，在这一理论中，我们处处把类函数式的对象——自守形式——替换为代数簇上的层；这使我们能够运用代数几何的有力方法来构造「自守层」。

编者感谢六位作者，他们以相当的功力把这些各异的部分缀合起来。我们还要感谢 C.~J.~Mozzochi 提供本书封面所用朗兰兹照片，感谢 Shlomit Davidzon 提供其他图形元素。最后，我们感谢耶路撒冷希伯来大学高等研究院——尤其是 Dahlia Aviely、Smedar Danziger、Pnina Feldman、Shani Freiman 与 Alex Levitzki——使本次讲习班得以举办。

\vspace{1.2em}
\begin{flushright}
Joseph Bernstein\\
特拉维夫大学\\
以色列拉马特阿维夫\\[0.8em]
Stephen Gelbart\\
魏茨曼科学研究所\\
Nicki and J.~Ira Harris 讲席教授\\
以色列雷霍沃特\\[0.8em]
2002 年 12 月
\end{flushright}
